\pagestyle{empty}
\tight
\begin{tblr}{width=\textwidth,colspec={|Q[c,m,1.9cm]|X[l,m]|},hlines,vlines,hspan=minimal,row{1}={bg=headblue},rowsep=1pt,colsep=3pt}
\SetCell{c,m}\hfil\logo\hfil & \textbf{POLITEKNIK NEGERI MALANG}\par\textbf{JURUSAN TEKNOLOGI INFORMASI}\par\textbf{PROGRAM STUDI: D4 TEKNIK INFORMATIKA} \\
\SetCell[c=2]{c} \textbf{RENCANA PEMBELAJARAN SEMESTER (RPS)} & \\
\end{tblr}
\begin{longtblr}{width=\textwidth,colspec={|Q[l,m,1.9cm]|X[0.85,l,m]|X[1.45,l,m]|X[0.75,l,m]|X[1.15,l,m]|},hlines,vlines,hspan=minimal,rowsep=1pt,colsep=2pt,row{1,3}={bg=headgray,font=\bfseries}}
MATA KULIAH & KODE & BOBOT (SKS)/jam & SEMESTER & TGL. PENYUSUNAN \\
Cloud Computing & RTI256107 & 2 SKS/4 jam & 6 & 1 Juli 2025 \\
OTORISASI & Dosen Pengembang RPS & Koordinator RMK & \SetCell[c=2]{l} Ka PRODI &  \\
 & Vipkas Al Hadid Firdaus, S.T., M.T. &  & \SetCell[c=2]{l}{\vspace{8mm}\par Dr. Ely Setyo Astuti, S.T., M.T.} &  \\
\SetCell[r=12]{l} Capaian Pembelajaran (CP) & \SetCell[c=4]{l,bg=headgray,font=\bfseries} Capaian Pembelajaran Lulusan Program Studi (CPL-Prodi) &  &  &  \\
 & CPL05 & \SetCell[c=3]{l} Mampu menghasilkan solusi teknologi komputasi dan infrastruktur digital yang sesuai dengan kebutuhan organisasi dan pengguna. &  &  \\
 & CPL09 & \SetCell[c=3]{l} Mampu menjelaskan cara kerja sistem komputer dan menerapkan pendekatan komputasi untuk menyelesaikan masalah. &  &  \\
 & CPL10 & \SetCell[c=3]{l} Mampu menganalisis persoalan kompleks dengan wawasan multidisiplin untuk menghasilkan solusi di bidang informatika dan ilmu komputer. &  &  \\
 & \SetCell[c=4]{l,bg=headgray,font=\bfseries} Capaian Pembelajaran mata kuliah (CPL-MK) &  &  &  \\
 & CPMK0509 & \SetCell[c=3]{l} Mahasiswa mampu merancang dan mengimplementasikan layanan komputasi virtual yang efisien, scalable, dan berkelanjutan. &  &  \\
 & CPMK0905 & \SetCell[c=3]{l} Mahasiswa mampu menjelaskan konsep virtualisasi dan prinsip cloud computing serta model layanan cloud. &  &  \\
 & CPMK1003 & \SetCell[c=3]{l} Mahasiswa mampu menganalisis dan merancang solusi jaringan on-premise/cloud yang andal, aman, dan scalable. &  &  \\
 & \SetCell[c=4]{l,bg=headgray,font=\bfseries} Kemampuan akhir tiap tahapan belajar (Sub-CPMK) &  &  &  \\
 & SCPMK0509-05101 & \SetCell[c=3]{l} Mahasiswa mampu merancang arsitektur cloud yang scalable menggunakan IaaS, PaaS, dan SaaS pada platform cloud utama. &  &  \\
 & SCPMK0905-05102 & \SetCell[c=3]{l} Mahasiswa mampu menjelaskan model virtualisasi, containerisasi (Docker, Kubernetes), dan arsitektur microservices. &  &  \\
 & SCPMK1003-05103 & \SetCell[c=3]{l} Mahasiswa mampu mengevaluasi dan memilih solusi cloud berdasarkan keandalan, keamanan, skalabilitas, dan biaya. &  &  \\
\SetCell[r=2]{l} Penilaian & \SetCell[c=4]{l} *Nilai CPMK didapat dari agregasi nilai SUB-CPMK yang dibebankan &  &  &  \\
 & \SetCell[c=4]{l}{\tiny\begin{tblr}{width=\linewidth,colspec={|Q[c,m,0.1\linewidth]|Q[c,m,0.16\linewidth]|X[l,m]|X[l,m]|X[l,m]|X[l,m]|X[l,m]|X[l,m]|Q[c,m,0.08\linewidth]|},hlines,vlines,hspan=minimal,rowsep=1pt,colsep=0.7pt,row{1,2}={bg=headgray,font=\bfseries}}
Kode CPMK & Kode SCPMK & \SetCell[c=6]{c} Bobot per Bentuk Penilaian & & & & & & Total Bobot \\
 & & Tugas & Kuis 1 & UTS & Kuis 2 & PBL & UAS & \\
CPMK0509 & SCPMK0509-05101 & 14\%: arsitektur cloud & 0\% & 6\%: arsitektur cloud & 0\% & 11\%: deployment cloud & 12\%: evaluasi arsitektur & 43\% \\
CPMK0905 & SCPMK0905-05102 & 10\%: virtualisasi, containerisasi & 5\%: konsep cloud & 7\%: virtualisasi dan model & 0\% & 0\% & 0\% & 22\% \\
CPMK1003 & SCPMK1003-05103 & 6\%: solusi cloud & 0\% & 0\% & 5\%: evaluasi solusi & 11\%: deployment & 13\%: evaluasi arsitektur & 35\% \\
\SetCell[c=8]{r}\textbf{Total Per Penilaian} & & & & & & & & \textbf{100\%} \\
\end{tblr}} &  &  &  \\
Diskripsi Singkat Mata Kuliah & \SetCell[c=4]{l} Cloud Computing membangun kemampuan merancang, mengimplementasikan, dan mengevaluasi solusi cloud menggunakan IaaS/PaaS/SaaS, containerisasi Docker-Kubernetes, dan prinsip arsitektur cloud yang andal dan scalable. &  &  &  \\
Bahan Kajian / Materi Pembelajaran & \SetCell[c=4]{l}\begin{enumerate}\item Virtualisasi, containerisasi, dan model layanan cloud (IaaS/PaaS/SaaS)\item Docker dan Kubernetes: containerisasi dan orkestrasi\item Arsitektur cloud: microservices, VPC, keamanan, dan jaringan\item Optimasi biaya, governance, dan strategi multi-cloud/hybrid\end{enumerate} &  &  &  \\
\SetCell[r=2]{l} Pustaka & \SetCell[c=4]{l} \textbf{Utama:}\begin{enumerate}\item Velte, T. 2010. Cloud Computing: A Practical Approach. McGraw-Hill.\item Burns, B. 2019. Kubernetes: Up and Running. O'Reilly.\item Dokumentasi AWS/GCP/Azure resmi.\end{enumerate} &  &  &  \\
 & \SetCell[c=4]{l} \textbf{Pendukung:} Materi LMS, dokumentasi resmi platform cloud, dan jobsheet praktikum. &  &  &  \\
Media Pembelajaran & \SetCell[c=2]{l} \textbf{Software:} AWS/GCP/Azure console, Docker, Kubernetes, LMS. &  & \SetCell[c=2]{l} \textbf{Hardware:} Komputer/laptop, proyektor. &  \\
Nama Dosen Pengampu & \SetCell[c=4]{l} Tim Pengajar Program Studi D4 TI &  &  &  \\
Matakuliah Syarat & \SetCell[c=4]{l} - &  &  &  \\
\end{longtblr}
\begin{longtblr}{width=\textwidth,colspec={|Q[c,m,0.06\textwidth]|X[1.35,l,m]|X[1.08,l,m]|X[1.16,l,m]|X[0.7,l,m]|X[1.24,l,m]|X[1.06,l,m]|X[0.98,l,m]|Q[c,m,0.05\textwidth]|},hlines,vlines,hspan=minimal,rowhead=2,row{1,2}={bg=headgreen,font=\bfseries},rowsep=1pt,colsep=1.5pt}
Mgg.\par Ke & Kemampuan Akhir Yang Direncanakan (Sub-CP-MK) & Bahan kajian (Materi Pembelajaran) & Bentuk dan Metode Pembelajaran & Estimasi Waktu & Pengalaman Belajar Mahasiswa & Kriteria \& Bentuk Penilaian & Indikator Penilaian & Bobot Penilaian (\%) \\
(1) & (2) & (3) & (4) & (5) & (6) & (7) & (8) & (9) \\
1 & SCPMK0905-05102 & Pengantar Cloud Computing: model dan karakteristik. & Modalitas: Blended Learning. Bentuk: Luring/Daring. Metode: Case Method, studi kasus, praktik cloud, dan peer review. & 1 x 2 x 50' tatap muka; 1 x 2 x 50' tugas/praktik mandiri. & Mahasiswa mengerjakan aktivitas terarah untuk pengantar cloud computing dan menunjukkan evidence hasil belajar. & Bentuk penilaian: Pengantar Cloud Computing. Mengacu rubrik RTM. & Ketepatan konsep; kualitas implementasi; validasi dan dokumentasi. & 2 \\
2 & SCPMK0905-05102 & Virtualisasi: hypervisor, VM, dan container. & Modalitas: Blended Learning. Bentuk: Luring/Daring. Metode: Case Method, studi kasus, praktik cloud, dan peer review. & 1 x 2 x 50' tatap muka; 1 x 2 x 50' tugas/praktik mandiri. & Mahasiswa mengerjakan aktivitas terarah untuk virtualisasi dan menunjukkan evidence hasil belajar. & Bentuk penilaian: Virtualisasi. Mengacu rubrik RTM. & Ketepatan konsep; kualitas implementasi; validasi dan dokumentasi. & 2 \\
3 & SCPMK0509-05101 & IaaS: komputasi, storage, dan jaringan virtual. & Modalitas: Blended Learning. Bentuk: Luring/Daring. Metode: Case Method, studi kasus, praktik cloud, dan peer review. & 1 x 2 x 50' tatap muka; 1 x 2 x 50' tugas/praktik mandiri. & Mahasiswa mengerjakan aktivitas terarah untuk IaaS dan menunjukkan evidence hasil belajar. & Bentuk penilaian: IaaS. Mengacu rubrik RTM. & Ketepatan konsep; kualitas implementasi; validasi dan dokumentasi. & 3 \\
4 & SCPMK0905-05102 & Kuis 1 konsep dan model layanan cloud. & Modalitas: Blended Learning. Bentuk: Luring/Daring. Metode: Case Method, studi kasus, praktik cloud, dan peer review. & 1 x 2 x 50' tatap muka; 1 x 2 x 50' tugas/praktik mandiri. & Mahasiswa mengerjakan aktivitas terarah untuk kuis 1 konsep dan model layanan cloud dan menunjukkan evidence hasil belajar. & Bentuk penilaian: Kuis 1 konsep dan model layanan cloud. Mengacu rubrik RTM. & Ketepatan konsep; kualitas implementasi; validasi dan dokumentasi. & 5 \\
5 & SCPMK0509-05101 & PaaS: platform development dan managed services. & Modalitas: Blended Learning. Bentuk: Luring/Daring. Metode: Case Method, studi kasus, praktik cloud, dan peer review. & 1 x 2 x 50' tatap muka; 1 x 2 x 50' tugas/praktik mandiri. & Mahasiswa mengerjakan aktivitas terarah untuk PaaS dan menunjukkan evidence hasil belajar. & Bentuk penilaian: PaaS. Mengacu rubrik RTM. & Ketepatan konsep; kualitas implementasi; validasi dan dokumentasi. & 2 \\
6 & SCPMK0905-05102 & SaaS dan model bisnis cloud. & Modalitas: Blended Learning. Bentuk: Luring/Daring. Metode: Case Method, studi kasus, praktik cloud, dan peer review. & 1 x 2 x 50' tatap muka; 1 x 2 x 50' tugas/praktik mandiri. & Mahasiswa mengerjakan aktivitas terarah untuk SaaS dan model bisnis cloud dan menunjukkan evidence hasil belajar. & Bentuk penilaian: SaaS dan model bisnis cloud. Mengacu rubrik RTM. & Ketepatan konsep; kualitas implementasi; validasi dan dokumentasi. & 2 \\
7 & SCPMK0905-05102 & Docker: containerisasi dan image management. & Modalitas: Blended Learning. Bentuk: Luring/Daring. Metode: Case Method, studi kasus, praktik cloud, dan peer review. & 1 x 2 x 50' tatap muka; 1 x 2 x 50' tugas/praktik mandiri. & Mahasiswa mengerjakan aktivitas terarah untuk Docker dan menunjukkan evidence hasil belajar. & Bentuk penilaian: Docker. Mengacu rubrik RTM. & Ketepatan konsep; kualitas implementasi; validasi dan dokumentasi. & 3 \\
8 & SCPMK0905-05102, SCPMK0509-05101 & UTS virtualisasi dan model layanan cloud. & Modalitas: Blended Learning. Bentuk: Luring/Daring. Metode: Case Method, studi kasus, praktik cloud, dan peer review. & 1 x 2 x 50' tatap muka; 1 x 2 x 50' tugas/praktik mandiri. & Mahasiswa mengerjakan aktivitas terarah untuk UTS virtualisasi dan model layanan cloud dan menunjukkan evidence hasil belajar. & Bentuk penilaian: UTS virtualisasi dan model layanan cloud. Mengacu rubrik RTM. & Ketepatan konsep; kualitas implementasi; validasi dan dokumentasi. & 13 \\
9 & SCPMK0509-05101 & Kubernetes: orkestrasi container dan deployment. & Modalitas: Blended Learning. Bentuk: Luring/Daring. Metode: Case Method, studi kasus, praktik cloud, dan peer review. & 1 x 2 x 50' tatap muka; 1 x 2 x 50' tugas/praktik mandiri. & Mahasiswa mengerjakan aktivitas terarah untuk Kubernetes dan menunjukkan evidence hasil belajar. & Bentuk penilaian: Kubernetes. Mengacu rubrik RTM. & Ketepatan konsep; kualitas implementasi; validasi dan dokumentasi. & 2 \\
10 & SCPMK0509-05101 & Microservices: arsitektur dan komunikasi. & Modalitas: Blended Learning. Bentuk: Luring/Daring. Metode: Case Method, studi kasus, praktik cloud, dan peer review. & 1 x 2 x 50' tatap muka; 1 x 2 x 50' tugas/praktik mandiri. & Mahasiswa mengerjakan aktivitas terarah untuk microservices dan menunjukkan evidence hasil belajar. & Bentuk penilaian: Microservices. Mengacu rubrik RTM. & Ketepatan konsep; kualitas implementasi; validasi dan dokumentasi. & 2 \\
11 & SCPMK1003-05103 & Kuis 2 Docker, Kubernetes, microservices. & Modalitas: Blended Learning. Bentuk: Luring/Daring. Metode: Case Method, studi kasus, praktik cloud, dan peer review. & 1 x 2 x 50' tatap muka; 1 x 2 x 50' tugas/praktik mandiri. & Mahasiswa mengerjakan aktivitas terarah untuk kuis 2 Docker, Kubernetes, microservices dan menunjukkan evidence hasil belajar. & Bentuk penilaian: Kuis 2 Docker, Kubernetes, microservices. Mengacu rubrik RTM. & Ketepatan konsep; kualitas implementasi; validasi dan dokumentasi. & 5 \\
12 & SCPMK1003-05103 & Cloud security: IAM, enkripsi, kepatuhan. & Modalitas: Blended Learning. Bentuk: Luring/Daring. Metode: Case Method, studi kasus, praktik cloud, dan peer review. & 1 x 2 x 50' tatap muka; 1 x 2 x 50' tugas/praktik mandiri. & Mahasiswa mengerjakan aktivitas terarah untuk cloud security dan menunjukkan evidence hasil belajar. & Bentuk penilaian: Cloud Security. Mengacu rubrik RTM. & Ketepatan konsep; kualitas implementasi; validasi dan dokumentasi. & 3 \\
13 & SCPMK1003-05103 & Cloud networking: VPC, load balancer, CDN. & Modalitas: Blended Learning. Bentuk: Luring/Daring. Metode: Case Method, studi kasus, praktik cloud, dan peer review. & 1 x 2 x 50' tatap muka; 1 x 2 x 50' tugas/praktik mandiri. & Mahasiswa mengerjakan aktivitas terarah untuk cloud networking dan menunjukkan evidence hasil belajar. & Bentuk penilaian: Cloud Networking. Mengacu rubrik RTM. & Ketepatan konsep; kualitas implementasi; validasi dan dokumentasi. & 3 \\
14 & SCPMK1003-05103 & Cost optimization dan cloud governance. & Modalitas: Blended Learning. Bentuk: Luring/Daring. Metode: Case Method, studi kasus, praktik cloud, dan peer review. & 1 x 2 x 50' tatap muka; 1 x 2 x 50' tugas/praktik mandiri. & Mahasiswa mengerjakan aktivitas terarah untuk cost optimization dan cloud governance dan menunjukkan evidence hasil belajar. & Bentuk penilaian: Cost Optimization dan Cloud Governance. Mengacu rubrik RTM. & Ketepatan konsep; kualitas implementasi; validasi dan dokumentasi. & 3 \\
15 & SCPMK1003-05103 & Multi-cloud dan hybrid cloud strategy. & Modalitas: Blended Learning. Bentuk: Luring/Daring. Metode: Case Method, studi kasus, praktik cloud, dan peer review. & 1 x 2 x 50' tatap muka; 1 x 2 x 50' tugas/praktik mandiri. & Mahasiswa mengerjakan aktivitas terarah untuk multi-cloud dan hybrid cloud strategy dan menunjukkan evidence hasil belajar. & Bentuk penilaian: Multi-cloud dan Hybrid Cloud Strategy. Mengacu rubrik RTM. & Ketepatan konsep; kualitas implementasi; validasi dan dokumentasi. & 3 \\
16 & SCPMK0509-05101, SCPMK1003-05103 & PBL deployment aplikasi di cloud. & Modalitas: Blended Learning. Bentuk: Luring/Daring. Metode: Case Method, studi kasus, praktik cloud, dan peer review. & 1 x 2 x 50' tatap muka; 1 x 2 x 50' tugas/praktik mandiri. & Mahasiswa mengerjakan aktivitas terarah untuk PBL deployment aplikasi di cloud dan menunjukkan evidence hasil belajar. & Bentuk penilaian: PBL deployment aplikasi di cloud. Mengacu rubrik RTM. & Ketepatan konsep; kualitas implementasi; validasi dan dokumentasi. & 22 \\
17 & SCPMK0509-05101, SCPMK1003-05103 & UAS evaluasi arsitektur cloud dan presentasi. & Modalitas: Blended Learning. Bentuk: Luring/Daring. Metode: Case Method, studi kasus, praktik cloud, dan peer review. & 1 x 2 x 50' tatap muka; 1 x 2 x 50' tugas/praktik mandiri. & Mahasiswa mengerjakan aktivitas terarah untuk UAS evaluasi arsitektur cloud dan presentasi dan menunjukkan evidence hasil belajar. & Bentuk penilaian: UAS evaluasi arsitektur cloud dan presentasi. Mengacu rubrik RTM. & Ketepatan konsep; kualitas implementasi; validasi dan dokumentasi. & 25 \\
\end{longtblr}
